\documentclass[12pt,a4paper]{ctexart}
\usepackage[margin=2.2cm]{geometry}
\usepackage{amsmath,amssymb,booktabs,array,longtable,enumitem}
\usepackage{mathtools}
\usepackage{fancyhdr}
\usepackage{titlesec}
\usepackage{hyperref}
\setCJKmainfont{Noto Serif CJK SC}
\setCJKsansfont{Noto Sans CJK SC}
\setmonofont{Noto Sans Mono CJK SC}
\hypersetup{colorlinks=false,pdfborder={0 0 0}}
\pagestyle{fancy}
\fancyhf{}
\lhead{有限差分推导整理}
\rhead{Problem 1}
\cfoot{\thepage}
\setlength{\parindent}{2em}
\setlength{\parskip}{0.35em}
\setlength{\headheight}{15pt}
\titleformat{\section}{\Large\bfseries}{\thesection}{0.8em}{}
\titleformat{\subsection}{\large\bfseries}{\thesubsection}{0.6em}{}
\newcommand{\dd}{\mathrm{d}}
\newcommand{\D}{\Delta}
\newcommand{\N}{\nabla}
\newcommand{\E}{E}
\newcommand{\half}{\frac12}
\newcommand{\fh}[1]{f(x+#1h)}
\newcommand{\fmh}[1]{f(x-#1h)}
\newcommand{\fx}{f(x)}
\newcommand{\fp}{f'}
\newcommand{\fpp}{f''}
\newcommand{\fiii}{f^{(3)}}
\newcommand{\fiv}{f^{(4)}}
\newcommand{\fv}{f^{(5)}}
\newcommand{\fvi}{f^{(6)}}
\newcommand{\fvii}{f^{(7)}}
\renewcommand{\arraystretch}{1.25}

\title{Problem 1：有限差分公式推导整理}
\author{}
\date{}

\begin{document}
\maketitle

\section*{题目}
推导 $y=f(x)$ 的二阶、三阶、四阶、五阶的向前差分、向后差分以及中心差分，并写成对应的导数公式。

\section{记号与差分算子}
设步长为 $h$，并记
\[
  x_i=x+ih,\qquad y_i=f(x_i).
\]
引入平移算子 $E$：
\[
  E f(x)=f(x+h),\qquad E^{-1}f(x)=f(x-h).
\]
因此
\[
  E^k f(x)=f(x+kh).
\]

向前差分、向后差分和中心差分分别定义为
\[
  \Delta=E-1,\qquad \nabla=1-E^{-1},\qquad \delta=E^{1/2}-E^{-1/2}.
\]
其中中心差分采用关于 $x$ 对称的点。偶数阶中心差分展开后已经是整数节点；奇数阶中心差分先得到半整数节点，再用平均算子
\[
  \mu=\frac{E^{1/2}+E^{-1/2}}{2}
\]
将其化为关于 $x$ 对称的整数节点形式。

\section{差分系数表}
\subsection{向前差分系数}
由
\[
  \Delta^n=(E-1)^n
\]
可得：
\[
\begin{array}{c|rrrrrr}
 n & E^0&E^1&E^2&E^3&E^4&E^5\\
\hline
1&-1&1&0&0&0&0\\
2&1&-2&1&0&0&0\\
3&-1&3&-3&1&0&0\\
4&1&-4&6&-4&1&0\\
5&-1&5&-10&10&-5&1
\end{array}
\]
例如
\[
  \Delta^5 f(x)=f(x+5h)-5f(x+4h)+10f(x+3h)-10f(x+2h)+5f(x+h)-f(x).
\]

\subsection{向后差分系数}
由
\[
  \nabla^n=(1-E^{-1})^n
\]
可得：
\[
\begin{array}{c|rrrrrr}
 n & E^0&E^{-1}&E^{-2}&E^{-3}&E^{-4}&E^{-5}\\
\hline
1&1&-1&0&0&0&0\\
2&1&-2&1&0&0&0\\
3&1&-3&3&-1&0&0\\
4&1&-4&6&-4&1&0\\
5&1&-5&10&-10&5&-1
\end{array}
\]

\subsection{中心差分系数与半整数节点的处理}
中心差分由
\[
  \delta^n=(E^{1/2}-E^{-1/2})^n
\]
给出。先列出原始系数表：
{\small
\[
\begin{array}{c|rrrrrrrrrrr}
 n&E^{-5/2}&E^{-2}&E^{-3/2}&E^{-1}&E^{-1/2}&1&E^{1/2}&E&E^{3/2}&E^2&E^{5/2}\\
\hline
1&0&0&0&0&-1&0&1&0&0&0&0\\
2&0&0&0&1&0&-2&0&1&0&0&0\\
3&0&0&-1&0&3&0&-3&0&1&0&0\\
4&0&1&0&-4&0&6&0&-4&0&1&0\\
5&-1&0&5&0&-10&0&10&0&-5&0&1
\end{array}
\]
}

实际写成导数公式时，要把半整数节点化为整数节点：

\begin{enumerate}[label=(\arabic*)]
\item 偶数阶中心差分已经是整数节点形式：
\[
  \delta^2=E-2+E^{-1},
\]
\[
  \delta^4=E^2-4E+6-4E^{-1}+E^{-2}.
\]

\item 奇数阶中心差分先乘平均算子 $\mu$，即
\[
  \mu=\frac{E^{1/2}+E^{-1/2}}{2}.
\]
三阶时：
\[
\begin{aligned}
\mu\delta^3
&=\frac{E^{1/2}+E^{-1/2}}{2}
  \left(E^{3/2}-3E^{1/2}+3E^{-1/2}-E^{-3/2}\right)\\
&=\frac12\left(E^2-2E+2E^{-1}-E^{-2}\right).
\end{aligned}
\]
所以三阶中心差分写成
\[
  \frac{f(x+2h)-2f(x+h)+2f(x-h)-f(x-2h)}{2h^3}.
\]
五阶时：
\[
\begin{aligned}
\mu\delta^5
&=\frac12\left(E^3-4E^2+5E-5E^{-1}+4E^{-2}-E^{-3}\right).
\end{aligned}
\]
所以五阶中心差分写成
\[
  \frac{f(x+3h)-4f(x+2h)+5f(x+h)-5f(x-h)+4f(x-2h)-f(x-3h)}{2h^5}.
\]
\end{enumerate}

\section{泰勒展开}
以下所有右端的 $f,f',f'',\ldots$ 均在 $x$ 处取值，即
\[
  f=f(x),\quad f'=f'(x),\quad f''=f''(x),\quad \ldots
\]
一般地，
\[
  f(x+kh)=\sum_{r=0}^{6}\frac{(kh)^r}{r!}f^{(r)}(x)+O(h^7).
\]
因此向前展开为
\[
\begin{aligned}
f(x+h)&=f+hf'+\frac{h^2}{2}f''+\frac{h^3}{6}f^{(3)}+\frac{h^4}{24}f^{(4)}+\frac{h^5}{120}f^{(5)}+\frac{h^6}{720}f^{(6)}+O(h^7),\\
f(x+2h)&=f+2hf'+2h^2f''+\frac43h^3f^{(3)}+\frac23h^4f^{(4)}+\frac4{15}h^5f^{(5)}+\frac4{45}h^6f^{(6)}+O(h^7),\\
f(x+3h)&=f+3hf'+\frac92h^2f''+\frac92h^3f^{(3)}+\frac{27}{8}h^4f^{(4)}+\frac{81}{40}h^5f^{(5)}+\frac{81}{80}h^6f^{(6)}+O(h^7),\\
f(x+4h)&=f+4hf'+8h^2f''+\frac{32}{3}h^3f^{(3)}+\frac{32}{3}h^4f^{(4)}+\frac{128}{15}h^5f^{(5)}+\frac{256}{45}h^6f^{(6)}+O(h^7),\\
f(x+5h)&=f+5hf'+\frac{25}{2}h^2f''+\frac{125}{6}h^3f^{(3)}+\frac{625}{24}h^4f^{(4)}+\frac{625}{24}h^5f^{(5)}+\frac{3125}{144}h^6f^{(6)}+O(h^7).
\end{aligned}
\]

向后推导时使用
\[
  f(x-kh)=\sum_{r=0}^{6}\frac{(-kh)^r}{r!}f^{(r)}(x)+O(h^7).
\]
也就是奇数阶导数项变号，偶数阶导数项不变。

中心差分求最低阶误差时，二阶、四阶需要展开到六阶即可；三阶、五阶的下一项含 $f^{(5)}$ 与 $f^{(7)}$，特别是五阶中心差分为了得到最低阶误差，需要补充到七阶：
\[
  f(x+kh)=\sum_{r=0}^{7}\frac{(kh)^r}{r!}f^{(r)}(x)+O(h^8).
\]

\section{向前差分推导及最低阶误差}
\subsection{二阶向前差分}
\[
\begin{aligned}
\Delta^2 f(x)
&=f(x+2h)-2f(x+h)+f(x)\\
&=h^2f''+h^3f^{(3)}+\frac7{12}h^4f^{(4)}+O(h^5).
\end{aligned}
\]
所以
\[
  \frac{\Delta^2 f(x)}{h^2}=f''+h f^{(3)}+O(h^2),
\]
即
\[
  \boxed{f''(x)=\frac{f(x+2h)-2f(x+h)+f(x)}{h^2}-h f^{(3)}(x)+O(h^2).}
\]
若只写差分近似，则
\[
  f''(x)\approx\frac{f(x+2h)-2f(x+h)+f(x)}{h^2},
\]
其最低阶误差为 $h f^{(3)}(x)$，精度为 $O(h)$。

\subsection{三阶向前差分}
\[
\begin{aligned}
\Delta^3 f(x)
&=f(x+3h)-3f(x+2h)+3f(x+h)-f(x)\\
&=h^3f^{(3)}+\frac32h^4f^{(4)}+O(h^5).
\end{aligned}
\]
因此
\[
  \boxed{f^{(3)}(x)=\frac{f(x+3h)-3f(x+2h)+3f(x+h)-f(x)}{h^3}-\frac32h f^{(4)}(x)+O(h^2).}
\]

\subsection{四阶向前差分}
\[
\begin{aligned}
\Delta^4 f(x)
&=f(x+4h)-4f(x+3h)+6f(x+2h)-4f(x+h)+f(x)\\
&=h^4f^{(4)}+2h^5f^{(5)}+O(h^6).
\end{aligned}
\]
因此
\[
  \boxed{f^{(4)}(x)=\frac{f(x+4h)-4f(x+3h)+6f(x+2h)-4f(x+h)+f(x)}{h^4}-2h f^{(5)}(x)+O(h^2).}
\]

\subsection{五阶向前差分}
\[
\begin{aligned}
\Delta^5 f(x)
&=f(x+5h)-5f(x+4h)+10f(x+3h)\\
&\quad -10f(x+2h)+5f(x+h)-f(x).
\end{aligned}
\]
用上面的泰勒展开逐项合并：
\[
\begin{array}{c|c}
\text{项} & \text{系数和}\\
\hline
f & 1-5+10-10+5-1=0\\
hf' & 5-5\cdot4+10\cdot3-10\cdot2+5=0\\
h^2f'' & \frac{25}{2}-5\cdot8+10\cdot\frac92-10\cdot2+5\cdot\frac12=0\\
h^3f^{(3)} & \frac{125}{6}-5\cdot\frac{32}{3}+10\cdot\frac92-10\cdot\frac43+5\cdot\frac16=0\\
h^4f^{(4)} & \frac{625}{24}-5\cdot\frac{32}{3}+10\cdot\frac{27}{8}-10\cdot\frac23+5\cdot\frac1{24}=0\\
h^5f^{(5)} & \frac{625}{24}-5\cdot\frac{128}{15}+10\cdot\frac{81}{40}-10\cdot\frac4{15}+5\cdot\frac1{120}=1\\
h^6f^{(6)} & \frac{3125}{144}-5\cdot\frac{256}{45}+10\cdot\frac{81}{80}-10\cdot\frac4{45}+5\cdot\frac1{720}=\frac52
\end{array}
\]
所以
\[
\begin{aligned}
\Delta^5 f(x)
&=h^5f^{(5)}+\frac52h^6f^{(6)}+O(h^7).
\end{aligned}
\]
因此
\[
  \boxed{f^{(5)}(x)=\frac{\Delta^5 f(x)}{h^5}-\frac52h f^{(6)}(x)+O(h^2).}
\]

\section{向后差分推导及最低阶误差}
向后差分的推导与向前差分类似，只是泰勒展开中奇数次幂项符号相反。

\subsection{二阶向后差分}
\[
\begin{aligned}
\nabla^2 f(x)
&=f(x)-2f(x-h)+f(x-2h)\\
&=h^2f''-h^3f^{(3)}+\frac7{12}h^4f^{(4)}+O(h^5).
\end{aligned}
\]
所以
\[
  \boxed{f''(x)=\frac{f(x)-2f(x-h)+f(x-2h)}{h^2}+h f^{(3)}(x)+O(h^2).}
\]

\subsection{三阶向后差分}
\[
\begin{aligned}
\nabla^3 f(x)
&=f(x)-3f(x-h)+3f(x-2h)-f(x-3h)\\
&=h^3f^{(3)}-\frac32h^4f^{(4)}+O(h^5).
\end{aligned}
\]
所以
\[
  \boxed{f^{(3)}(x)=\frac{f(x)-3f(x-h)+3f(x-2h)-f(x-3h)}{h^3}+\frac32h f^{(4)}(x)+O(h^2).}
\]

\subsection{四阶向后差分}
\[
\begin{aligned}
\nabla^4 f(x)
&=f(x)-4f(x-h)+6f(x-2h)-4f(x-3h)+f(x-4h)\\
&=h^4f^{(4)}-2h^5f^{(5)}+O(h^6).
\end{aligned}
\]
所以
\[
  \boxed{f^{(4)}(x)=\frac{f(x)-4f(x-h)+6f(x-2h)-4f(x-3h)+f(x-4h)}{h^4}+2h f^{(5)}(x)+O(h^2).}
\]

\subsection{五阶向后差分}
\[
\begin{aligned}
\nabla^5 f(x)
&=f(x)-5f(x-h)+10f(x-2h)-10f(x-3h)\\
&\quad +5f(x-4h)-f(x-5h)\\
&=h^5f^{(5)}-\frac52h^6f^{(6)}+O(h^7).
\end{aligned}
\]
所以
\[
  \boxed{f^{(5)}(x)=\frac{\nabla^5 f(x)}{h^5}+\frac52h f^{(6)}(x)+O(h^2).}
\]

\section{中心差分推导及最低阶误差}
\subsection{二阶中心差分}
\[
\begin{aligned}
C_2&=f(x+h)-2f(x)+f(x-h)\\
&=h^2f''+\frac1{12}h^4f^{(4)}+\frac1{360}h^6f^{(6)}+O(h^8).
\end{aligned}
\]
故
\[
  \boxed{f''(x)=\frac{f(x+h)-2f(x)+f(x-h)}{h^2}-\frac{h^2}{12}f^{(4)}(x)+O(h^4).}
\]

\subsection{三阶中心差分}
三阶中心差分用整数节点写为
\[
  C_3=f(x+2h)-2f(x+h)+2f(x-h)-f(x-2h).
\]
由泰勒展开合并得
\[
  C_3=2h^3f^{(3)}+\frac12h^5f^{(5)}+O(h^7).
\]
所以
\[
  \frac{C_3}{2h^3}=f^{(3)}+\frac14h^2f^{(5)}+O(h^4),
\]
即
\[
  \boxed{f^{(3)}(x)=\frac{f(x+2h)-2f(x+h)+2f(x-h)-f(x-2h)}{2h^3}-\frac{h^2}{4}f^{(5)}(x)+O(h^4).}
\]

\subsection{四阶中心差分}
\[
\begin{aligned}
C_4&=f(x+2h)-4f(x+h)+6f(x)-4f(x-h)+f(x-2h)\\
&=h^4f^{(4)}+\frac16h^6f^{(6)}+O(h^8).
\end{aligned}
\]
所以
\[
  \boxed{f^{(4)}(x)=\frac{f(x+2h)-4f(x+h)+6f(x)-4f(x-h)+f(x-2h)}{h^4}-\frac{h^2}{6}f^{(6)}(x)+O(h^4).}
\]

\subsection{五阶中心差分}
五阶中心差分用整数节点写为
\[
  C_5=f(x+3h)-4f(x+2h)+5f(x+h)-5f(x-h)+4f(x-2h)-f(x-3h).
\]
为得到最低阶误差，展开到七阶：
\[
  C_5=2h^5f^{(5)}+\frac23h^7f^{(7)}+O(h^9).
\]
于是
\[
  \frac{C_5}{2h^5}=f^{(5)}+\frac13h^2f^{(7)}+O(h^4),
\]
即
\[
\boxed{\begin{aligned}
f^{(5)}(x)&=\frac{f(x+3h)-4f(x+2h)+5f(x+h)-5f(x-h)+4f(x-2h)-f(x-3h)}{2h^5}\\
&\quad -\frac{h^2}{3}f^{(7)}(x)+O(h^4).
\end{aligned}}
\]

\section{最终公式汇总}
下面按阶数汇总，先给出主要差分商，再注明最低阶误差。

\subsection*{二阶导数}
\[
\begin{aligned}
f''(x)&=\frac{f(x+2h)-2f(x+h)+f(x)}{h^2}-h f^{(3)}(x)+O(h^2),\\
f''(x)&=\frac{f(x)-2f(x-h)+f(x-2h)}{h^2}+h f^{(3)}(x)+O(h^2),\\
f''(x)&=\frac{f(x+h)-2f(x)+f(x-h)}{h^2}-\frac{h^2}{12}f^{(4)}(x)+O(h^4).
\end{aligned}
\]

\subsection*{三阶导数}
\[
\begin{aligned}
f^{(3)}(x)&=\frac{f(x+3h)-3f(x+2h)+3f(x+h)-f(x)}{h^3}-\frac32h f^{(4)}(x)+O(h^2),\\
f^{(3)}(x)&=\frac{f(x)-3f(x-h)+3f(x-2h)-f(x-3h)}{h^3}+\frac32h f^{(4)}(x)+O(h^2),\\
f^{(3)}(x)&=\frac{f(x+2h)-2f(x+h)+2f(x-h)-f(x-2h)}{2h^3}-\frac{h^2}{4}f^{(5)}(x)+O(h^4).
\end{aligned}
\]

\subsection*{四阶导数}
\[
\begin{aligned}
f^{(4)}(x)&=\frac{f(x+4h)-4f(x+3h)+6f(x+2h)-4f(x+h)+f(x)}{h^4}-2h f^{(5)}(x)+O(h^2),\\
f^{(4)}(x)&=\frac{f(x)-4f(x-h)+6f(x-2h)-4f(x-3h)+f(x-4h)}{h^4}+2h f^{(5)}(x)+O(h^2),\\
f^{(4)}(x)&=\frac{f(x+2h)-4f(x+h)+6f(x)-4f(x-h)+f(x-2h)}{h^4}-\frac{h^2}{6}f^{(6)}(x)+O(h^4).
\end{aligned}
\]

\subsection*{五阶导数}
\[
\begin{aligned}
f^{(5)}(x)&=\frac{f(x+5h)-5f(x+4h)+10f(x+3h)-10f(x+2h)+5f(x+h)-f(x)}{h^5}\\
&\quad -\frac52h f^{(6)}(x)+O(h^2),\\[0.4em]
f^{(5)}(x)&=\frac{f(x)-5f(x-h)+10f(x-2h)-10f(x-3h)+5f(x-4h)-f(x-5h)}{h^5}\\
&\quad +\frac52h f^{(6)}(x)+O(h^2),\\[0.4em]
f^{(5)}(x)&=\frac{f(x+3h)-4f(x+2h)+5f(x+h)-5f(x-h)+4f(x-2h)-f(x-3h)}{2h^5}\\
&\quad -\frac{h^2}{3}f^{(7)}(x)+O(h^4).
\end{aligned}
\]

\end{document}
